\title{Prompting Science Report 2: The Decreasing Value of Chain of Thought in Prompting}
\begin{document}
\maketitle

\begin{abstract}
Prompting Science Report 2: The Decreasing Value of Chain of Thought in Prompting Lennart Meincke1,2, Ethan Mollick1, Lilach Mollick1, Dan Shapiro1,3 1 Generative AI Labs, The Wharton School of Business, University of Pennsylvania​ 2 WHU–Otto Beisheim School of Management​ 3 Glowforge Summary This is the second in a series of short reports that seek to help business, education, and policy leaders understand the technical details of working with AI through rigorous testing. In this report, we investigate Chain-of-Thought (CoT) prompting, a technique that encourages a large language model (LLM) to "think step by step" (Wei et al., 2022). CoT is a widely adopted method for improving reasoning tasks, however, our findings reveal a more nuanced picture of its effectiveness. We demonstrate two things: ●​ The effectiveness of Chain-of-Thought prompting can vary greatly depending on the type of task and model. For non-reasoning models, CoT generally improves average performance by a small amount, particularly if the model does not inherently engage in step-by-step processing by default. However, CoT can introduce more variability in answers, sometimes triggering occasional errors in questions the model would otherwise get right. We also found that many recent models perform some form of CoT reasoning even if not asked (see Table S1); for these models, a request to perform CoT had little impact. Performing CoT generally requires far more tokens (increasing cost and time) than direct a
\end{abstract}

\section{Recovered Text}
Prompting Science Report 2: The Decreasing Value of Chain of Thought in
Prompting
Lennart Meincke1,2, Ethan Mollick1, Lilach Mollick1, Dan Shapiro1,3
1 Generative AI Labs, The Wharton School of Business, University of Pennsylvania​
2 WHU–Otto Beisheim School of Management​
3 Glowforge
Summary
This is the second in a series of short reports that seek to help business, education, and policy
leaders understand the technical details of working with AI through rigorous testing. In this
report, we investigate Chain-of-Thought (CoT) prompting, a technique that encourages a large
language model (LLM) to "think step by step" (Wei et al., 2022). CoT is a widely adopted
method for improving reasoning tasks, however, our findings reveal a more nuanced picture of
its effectiveness. We demonstrate two things:
●​ The effectiveness of Chain-of-Thought prompting can vary greatly depending on the type
of task and model. For non-reasoning models, CoT generally improves average
performance by a small amount, particularly if the model does not inherently engage in
step-by-step processing by default. However, CoT can introduce more variability in
answers, sometimes triggering occasional errors in questions the model would otherwise
get right. We also found that many recent models perform some form of CoT reasoning
even if not asked (see Table S1); for these models, a request to perform CoT had little
impact. Performing CoT generally requires far more tokens (increasing cost and time)
than direct answers.
●​ For models designed with explicit reasoning capabilities, CoT prompting often results in
only marginal, if any, gains in answer accuracy. However, it significantly increases the
time and tokens needed to generate a response.
Taken together, this suggests that a simple CoT prompt is generally still a useful tool for
boosting average performance in non-reasoning models, especially older or smaller models that
may not engage in a CoT reasoning by default. However, the gains must be weighed against
increased response times and potential decreases in perfect accuracy due to more variability in
answers. For dedicated reasoning models, the added benefits of explicit CoT prompting appear
negligible and may not justify the substantial increase in processing time.
How we Benchmark the AI
In this Report, our focus is on understanding the impact of Chain-of-Thought prompting across
different types of models. The goal of our report is not to assess model performance on any
given benchmark, but rather to show how the effectiveness of Chain-of-Thought prompting can
vary greatly depending on the type of task and model.
For a benchmark, we selected the commonly-used GPQA Diamond (Graduate-Level
Google-Proof Q\&A Benchmark) dataset (Rein et al. 2024). The GPQA Diamond set comprises
198 multiple-choice PhD-level questions across biology, physics, and chemistry. This is a
challenging test: PhDs in the corresponding domains reach 65\% accuracy (74\% when
discounting clear mistakes the experts identified in retrospect), while highly skilled non-expert
validators only reach 34\% accuracy, despite spending on average over 30 minutes with
unrestricted access to the web (i.e., the questions are "Google-proof)” (Rein et al. 2024).
We specifically look at common non-reasoning and reasoning models:
●​ Non-Reasoning Models: Sonnet 3.5 (claude-3-5-sonnet-20240620), Gemini 2.0 Flash
(gemini-2.0-flash-001), GPT-4o-mini (gpt-4o-mini-2024-07-18), GPT-4o
(gpt-4o-2024-08-06), and Gemini Pro 1.5 (gemini-1.5-pro-001).
●​ Reasoning Models: For this category, we examined models that engage in an initial
reasoning process: o3-mini (o3-mini-2025-01-31), o4-mini (o4-mini-2025-04-16), and
Flash 2.5 (gemini-2.5-flash-preview-05-20).
Each question in each condition is asked in 25 separate trials to ensure a more rigorous
analysis, since LLM responses can often vary (Miller 2024). In our tests, we find that the
correlation between all performance metrics on 25 trials and higher numbers of trials, such as
100, is very high (see Table S2), but this may be task dependent. We follow the rating measures
from our previous report (Meincke et al. 2025) and report the following metrics:
●​ 100\% Correct: In this condition, we require the AI to get the answer correct on all trials
(25/25) of a question to count that question as successful. This is our strictest condition
for tasks where there is no room for errors.
●​ 90\% Correct: We relax the 100\% condition and allow a few errors per question, but
23/25 answers must be correct. This might be comparable to human acceptable error
rates.
●​ 51\% Correct: A simple majority vote where AI needs to get at least 13/25 answers on a
question correct.
●​ Average Rating: We do not collapse measurements per-question, but instead use the
overall average of all trials across all questions (N = 4950).
We generally do not consider the more commonly used PASS@N (passing when one of N
answers is correct) and CONSENSUS (passing when the modal answer is correct) metrics
useful for benchmarking real-world applications.
How we Prompt the AI
We follow the standard implementation by Rein et al. (2024) for GPQA with the default system
prompt1  and temperature set to 0. At the time of writing, most non-reasoning models have a
tendency to produce a short reasoning trace before answering, similar to CoT. To isolate the
effectiveness of a CoT prompt, we specifically prompt the model to answer without thinking. In
total, we compare the following prompts (see Table S1 for example conversations):
●​ Answer directly prompt (“Direct"): We instruct the LLM to answer without generating
additional reasoning tokens (“Answer directly without any explanation or thinking. Just
provide the answer.”)
●​ Think step by step (“Step by step”): In this variation, we provide the LLM with a
simple CoT prompt that instructs it to “think” step by step before providing an answer
(“Think step by step”).  While this is a simple version of CoT, in our tests, different CoT
prompt variants had negligible effects (see Figure S3).
●​ No prompt (“Default”): We provide no specific suffix to the question and also remove
the formatting constraint (“Format your response as follows: ‘The correct answer is
(insert answer here)’”), allowing the model to choose how to best reply to the question. In
most cases, this produces a short CoT-like thinking output before answering the
question.
Each prompt condition was tested 25 times each across all 198 questions of the Diamond
GPQA dataset (4,950 runs per prompt per model).
Results
Across all non-reasoning models tested, employing a "Step-by-Step" (CoT) prompt generally led
to an improvement in the "Average rating" compared to the "Direct" prompt across all models
(see Figure 1). The effect was significant and most pronounced for Gemini Flash 2.0 (RD =
0.135, p < .001) and Sonnet 3.5 (RD = 0.117, p < .001) with GPT 4o-mini showing the smallest
and insignificant gains (RD = 0.044, p = .067). However, the impact on how many questions met
the "100\% Correct" (25/25 correct) metric was more varied. While Sonnet 3.5 saw a large
significant increase (RD = 0.101, p = .001), GPT-4o’s saw no significant improvement (RD =
0.01, p = .624). All other models saw significant declines, which were especially sharp for
Gemini Flash 2.0 (RD = -0.131, p < .001) and Gemini Pro 1.5 (RD = -0.172, p < .001). Using the
1 “You are a very intelligent assistant, who follows instructions directly.”
90\% metric(at least 23/25 correct), Sonnet 3.5 similarly saw significant gains (RD = 0.136, p <
.001), whereas GPT-4o-mini and Gemini Pro 1.5 saw significant losses (RD = -0.066, p = .033;
RD = -0.131, p < .001). All models gained in the majority correctness condition (at least 13/25
correct), with significant effects for Sonnet 3.5 (RD = 0.131, p < .001), Gemini Flash 2.0 (RD =
0.136, p < .001) and GPT-4o (RD = 0.071, p = 0.042). CoT requests took between 35-600\%
(5-15 seconds) longer than direct requests (see Figure S1). Confidence intervals for
comparisons are reported in Table S3.
Figure 1. GPQA Diamond performance across non-reasoning models comparing an immediate
answer with Chain-of-Thought.
Note. The error bars show 95\% confidence intervals for individual proportions. For detailed
statistical comparisons between conditions, see Table S3.
Comparing the performance of the CoT intervention to the default model behavior without any
specific prompt or formatting constraints, more closely comparable to a standard
consumer-facing chat, produced mixed results. Gains from CoT prompting were modest.
Investigating individual model replies showed that the model would often perform CoT by default
even without explicitly prompting it to think step by step. Thus, the value of explicitly prompting
for CoT greatly diminished and led to much smaller improvements in performance. There was a
statistically significant average increase for Gemini Flash 2.0 (RD = 0.062, p < .001) and
GPT-4o (RD = 0.069, p = 0.003) but no significant change for Sonnet 3.5 (RD = -0.019, p =
0.189) and GPT-4o-mini (RD = 0.004, p = .728 ;see Figure 2).
When using the "100\% Correct" metric, Gemini Flash 2.0 saw significant gains (RD = 0.071, p =
.039). Sonnet 3.5 saw small gains (RD = 0.020, p = .456) and GPT-4o (RD = -0.010, p = .639)
and GPT-4o-mini (RD = -0.005, p = .730) small decreases, but all were insignificant. For the
"90\% Correct" metric, Gemini Flash 2.0 once again improved significantly (RD = 0.111, p <
.001). Results for all other models were insignificant. Using the 51\% metric, no model showed
significant changes. Confidence intervals for comparisons are reported in Table S4.
Figure 2. GPQA Diamond performance across non-reasoning models comparing an
unprompted answer with Chain-of-Thought.
Note. The error bars show 95\% confidence intervals for individual proportions. For detailed
statistical comparisons between conditions, see Table S4.
For reasoning models, Chain-of-Thought prompting yielded only marginal improvements in
accuracy (see Figure 3). Both o3 and o4 slightly and significantly benefited from CoT on
average (RD = 0.029, p = .024; RD = 0.031, p = .003), whereas Gemini Flash 2.5’s performance
saw a significant small decrease (RD = -0.033, p = .005). Using the threshold metrics, we found
almost no significant changes. The only significant effects were a small gain for o4-mini at 51\%
(RD = 0.056, p = .003) and a performance decrease for Gemini Flash 2.5 at 100\% and 90\% (RD
= -0.131, p < .001; RD = -0.071, p = .016). CoT requests took between 20-80\% (10-20 seconds)
longer than direct requests (see Figure S2). Confidence intervals for comparisons are reported
in Table S5.
Figure 3. GPQA Diamond performance across reasoning models comparing an immediate
answer with Chain-of-Thought.
Note. The error bars show 95\% confidence intervals for individual proportions. For detailed
statistical comparisons between conditions, see Table S5.
Discussion and Conclusion
Our findings indicate that Chain-of-Thought prompting is not a universally optimal strategy.
For non-reasoning models, a simple CoT prompt can be a valuable tool for enhancing average
performance, particularly when the model doesn't already engage in step-by-step processing. In
three of the five non-reasoning models, CoT introduced a higher likelihood of errors on "easy"
questions that the model would otherwise get right, harming performance on the "100\% Correct"
metric. Performance on other, more challenging questions increased, however, leading to higher
average performance when using CoT. Chain-of-thought increases the token count, time, and
cost required to answer a prompt.
We found that many non-reasoning models perform a version of Chain-of-Thought even if
unprompted, which makes a CoT prompt less helpful. Additionally, we found the common
practice of prompting a model to reply with only the answer and nothing else is likely to harm the
performance of non-reasoning models, as it prevents them from performing CoT reasoning.
For reasoning models, the benefits of explicit CoT prompting on answer accuracy are marginal
at best. Given the substantial increase in response time and token usage, the utility of applying
external CoT prompts to these models is questionable for most practical purposes.
This study has several limitations, including the fact that we tested a limited set of models, only
one benchmark and that we used a simple version of Chain-of-Thought in our tests. However,
the fact that our tests of more sophisticated CoT prompts did not yield meaningful differences in
outputs, and that models were fairly consistent in their responses to CoT, do provide some
evidence that these results are general in nature. Highly customized CoT prompts aimed at a
specific problem may yield larger benefits, but generic approaches to CoT have more limited
results. Ultimately, the decision to use CoT prompting should be guided by the specific model's
characteristics, the task requirements, and a careful evaluation of the balance between desired
accuracy improvements, tolerance for variability, and acceptable response latency.
Acknowledgments
OpenAI provided API credits used in this research. No party other than the authors participated
in the design, execution, analysis, or interpretation of this study; and the authors receive no
financial compensation from the companies in this study, including OpenAI, Google, and
Anthropic.
References
Meincke L, Mollick ER, Mollick L, Shapiro D (2025) Prompting Science Report 1: Prompt
Engineering is Complicated and Contingent. (March 4)
https://papers.ssrn.com/abstract=5165270.
Miller E (2024) Adding Error Bars to Evals: A Statistical Approach to Language Model
Evaluations. (November 1) http://arxiv.org/abs/2411.00640.
Rein D, Hou BL, Stickland AC, Petty J, Pang RY, Dirani J, Michael J, Bowman SR (2024)
GPQA: A Graduate-Level Google-Proof Q\&A Benchmark. First Conference on
Language Modeling.
Wei J, Wang X, Schuurmans D, Bosma M, Ichter B, Xia F, Chi EH, Le QV, Zhou D (2022)
Chain-of-thought prompting elicits reasoning in large language models. Proceedings of
the 36th International Conference on Neural Information Processing Systems. NIPS
’22. (Curran Associates Inc., Red Hook, NY, USA), 24824–24837.
Supplementary Information
Figures
Figure S1. Average time per question in seconds for non-reasoning models on GPQA Diamond.
Error bars represent the 95\% confidence interval for each individual estimate. N = 197 per
condition.
Figure S2. Average time per question in seconds for reasoning models on GPQA Diamond.
Error bars represent the 95\% confidence interval for each individual estimate. N = 197 per
condition (196 for Flash 2.5 as one run took  much longer due to rate-limits).
Figure S3. Comparison of different Chain-of-Thought prompt variants on GPQA Diamond using
GPT-4o at 100 samples per question across four conditions: Step by Step (“Think step by
step”), Think Carefully (“Think very carefully before you provide your response. Consider
multiple different perspectives first. Then, weigh the pros and cons of each perspective. Finally,
provide your answer.), Decompose Problem (“First, decompose the problem. Then, consider
viable paths to solve it. Then, weigh the paths and pick the solution.) and Weigh Paths (“Before
you answer, carefully weigh three different possible solution paths. Explore them deeply. Then,
make your final answer based on your reasoning.)
Tables
Table S1. Example GPQA questions and comparison of answers per condition from GPT-4o.
Answer Directly
Think step by step
No prompt + no formatting
constraint
What is the correct answer to
this question: If a sperm from
species A is injected into an
egg from species B and both
species have the same
number of chromosomes,
what would be the main
cause of the resulting zygote
mortality?
A) Species specific zona
pellucida proteins on the egg
cannot bind sperms from a
different species.
B) Epistatic interactions
between the genes of
different species
C) Chromosomal
incompatibilities will cause
failure of meiosis leading to
death of zygote.
D) Chromosomal
recombination will not occur
in different species.
Format your response as
follows: "The correct answer
is (insert answer here)"
Answer directly without any
explanation or thinking. Just
provide the answer.
What is the correct answer to
this question: If a sperm from
species A is injected into an
egg from species B and both
species have the same
number of chromosomes,
what would be the main
cause of the resulting zygote
mortality?
A) Species specific zona
pellucida proteins on the egg
cannot bind sperms from a
different species.
B) Epistatic interactions
between the genes of
different species
C) Chromosomal
incompatibilities will cause
failure of meiosis leading to
death of zygote.
D) Chromosomal
recombination will not occur
in different species.
Format your response as
follows: "The correct answer
is (insert answer here)" Think
step by step.
What is the correct answer to
this question: If a sperm from
species A is injected into an
egg from species B and both
species have the same
number of chromosomes,
what would be the main
cause of the resulting zygote
mortality?
A) Species specific zona
pellucida proteins on the egg
cannot bind sperms from a
different species.
B) Epistatic interactions
between the genes of
different species
C) Chromosomal
incompatibilities will cause
failure of meiosis leading to
death of zygote.
D) Chromosomal
recombination will not occur
in different species.
Answers
The correct answer is B)
Epistatic interactions between
the genes of different
species.
To determine the correct
answer, let's analyze each
option step by step:
A) Species specific zona
pellucida proteins on the egg
cannot bind sperms from a
different species.
- This option describes a
barrier to fertilization, where
the proteins on the egg's
The question describes a
scenario where fertilization
has already been technically
achieved by directly injecting
a sperm from species A into
an egg from species B,
bypassing the typical
biological barriers like
species-specific binding
(which would be relevant if
fertilization were occurring
zona pellucida would prevent
sperm from a different
species from binding and
fertilizing the egg. However,
the question states a sperm
from species A is already
injected into an egg from
species B, implying
fertilization has already
bypassed these barriers.
Therefore, this is not a cause
of post-fertilization zygote
mortality.
B) Epistatic interactions
between the genes of
different species.
- Epistasis refers to the
interaction of genes at
different loci, where the
expression of one gene is
affected by one or more other
genes. In a hybrid organism,
epistatic interactions between
incompatible genes from
different species could lead to
developmental issues and
cause zygote mortality. This
is a plausible cause of zygote
mortality after fertilization.
C) Chromosomal
incompatibilities will cause
failure of meiosis leading to
death of zygote.
- Meiosis is the process of
forming gametes (sperm and
eggs), and it does not occur
in a zygote. After fertilization,
the zygote undergoes mitotic
divisions, not meiotic ones.
Therefore, chromosomal
incompatibilities related to
meiosis are not relevant to
zygote mortality
post-fertilization.
Chromosomal
incompatibilities can affect
zygote development, but they
are not related to the failure
naturally). Therefore, option A
is not relevant to the scenario
after injection.
Once the sperm from species
A is inside the egg from
species B, the main issue
arises at the genetic and
chromosomal level. Even if
both species have the same
number of chromosomes, the
combination of genetic
material from two different
species can lead to problems.
Option B, "epistatic
interactions between the
genes of different species,"
suggests that there could be
incompatible interactions
between the genes coming
from each parent, which
might lead to developmental
problems.
However, one of the most
pressing issues is usually
chromosomal incompatibility,
which could lead to failures in
zygotic development. The
correct answer would likely
be option C: "Chromosomal
incompatibilities will cause
failure of meiosis leading to
death of zygote." Although
meiosis specifically is not
involved here (since we are
dealing with a zygote),
chromosomal
incompatibilities often
interfere with the normal cell
division and development
processes.
In summary, the most
accurate answer in terms of
explaining zygote mortality
due to interspecies
fertilization at the
chromosomal level would be
of meiosis in this context.
D) Chromosomal
recombination will not occur
in different species.
- Chromosomal
recombination occurs during
the formation of gametes
(meiosis), not during the
development of a zygote.
Post-fertilization mitotic
divisions do not involve
recombination. Additionally,
even if recombination issues
were relevant, the question
specifically asks about zygote
mortality, not about
recombination processes
before fertilization.
Considering the above
analysis, option B (epistatic
interactions between the
genes of different species) is
the most likely cause of
zygote mortality, as it could
lead to developmental issues
due to incompatible gene
interactions from two different
species.
The correct answer is B)
Epistatic interactions between
the genes of different
species.
C.
Table S2. Precision and power comparison for per-question sample sizes of 25 versus 100 trials
using GPT-4o mini in the direct condition. Point estimates and 95 \% percentile bootstrap
confidence intervals (5,000 replicates) were obtained by resampling each question’s latent
success probability at the indicated trial count. “Average” refers to the overall average rating
across all questions; “100\%,” “90\%”, “51\%” list the fraction of questions whose simulated rating
meets or exceeds that threshold. A paired bootstrap test (5,000 permutations) confirmed that no
metric reached statistical significance at α=0.05, indicating that 25 trials per question provide
estimation precision and statistical power comparable to 100 trials for the effect sizes
considered in this study. Multiple trial runs with 25 and 100 samples each confirmed the power
simulation findings.
N Trials
Metric
Estimate [95\% CI]
25
Average
0.389 [0.382, 0.396]
100
Average
0.389 [0.385, 0.393]
25
100\%
0.185 [0.162, 0.207]
100
100\%
0.147 [0.136, 0.162]
25
51\%
0.380 [0.364, 0.394]
100
51\%
0.377 [0.369, 0.389]
25
90\%
0.247 [0.227, 0.263]
100
90\%
0.249 [0.237, 0.263]
Table S3. Pairwise-comparison results between direct and CoT responses for each model using
paired bootstrap-permutation tests (5,000 replicates). P-values represent the proportion of
permuted differences with absolute values exceeding observed differences under the null
hypothesis. Questions were considered 'perfect' at ≥100 successes, with sensitivity analyses at
thresholds of 90 and 51.
Model
Conditions
Comparison
RD [95\% CI]
Statistics
Sonnet 3.5
Direct - Step by Step
Average
0.117 [0.069, 0.166]
p < 0.001
Sonnet 3.5
Direct - Step by Step
100\%
0.101 [0.040, 0.162]
p = 0.001
Sonnet 3.5
Direct - Step by Step
90\%
0.136 [0.071, 0.202]
p < 0.001
Sonnet 3.5
Direct - Step by Step
51\%
0.131 [0.056, 0.207]
p = 0.001
Gemini 2.0 Flash
Direct - Step by Step
Average
0.135 [0.076, 0.199]
p < 0.001
Gemini 2.0 Flash
Direct - Step by Step
100\%
-0.131 [-0.207, -0.056]
p < 0.001
Gemini 2.0 Flash
Direct - Step by Step
90\%
-0.030 [-0.101, 0.045]
p = 0.381
Gemini 2.0 Flash
Direct - Step by Step
51\%
0.136 [0.061, 0.212]
p < 0.001
GPT-4o-mini
Direct - Step by Step
Average
0.044 [-0.003, 0.090]
p = 0.067
GPT-4o-mini
Direct - Step by Step
100\%
-0.091 [-0.141, -0.035]
p = 0.001
GPT-4o-mini
Direct - Step by Step
90\%
-0.066 [-0.126, -0.010]
p = 0.033
GPT-4o-mini
Direct - Step by Step
51\%
0.020 [-0.045, 0.086]
p = 0.609
GPT-4o
Direct - Step by Step
Average
0.072 [0.028, 0.118]
p = 0.003
GPT-4o
Direct - Step by Step
100\%
0.010 [-0.040, 0.061]
p = 0.624
GPT-4o
Direct - Step by Step
90\%
0.035 [-0.015, 0.086]
p = 0.150
GPT-4o
Direct - Step by Step
51\%
0.071 [0.000, 0.141]
p = 0.042
Gemini Pro 1.5
Direct - Step by Step
Average
0.071 [0.018, 0.128]
p = 0.014
Gemini Pro 1.5
Direct - Step by Step
100\%
-0.172 [-0.242, -0.101]
p < 0.001
Gemini Pro 1.5
Direct - Step by Step
90\%
-0.131 [-0.197, -0.066]
p < 0.001
Gemini Pro 1.5
Direct - Step by Step
51\%
0.056 [-0.015, 0.126]
p = 0.097
Table S4. Pairwise-comparison results between default and CoT responses for each
non-reasoning model using paired bootstrap-permutation tests (5,000 replicates). P-values
represent the proportion of permuted differences with absolute values exceeding observed
differences under the null hypothesis. Questions were considered 'perfect' at ≥100 successes,
with sensitivity analyses at thresholds of 90 and 51.
Model
Conditions
Metric
RD [95\% CI]
Statistics
Sonnet 3.5
Default - Step by Step
Average
-0.019 [-0.047, 0.009]
p = 0.189
Sonnet 3.5
Default - Step by Step
100\%
0.020 [-0.040, 0.081]
p = 0.456
Sonnet 3.5
Default - Step by Step
90\%
-0.020 [-0.076, 0.035]
p = 0.515
Sonnet 3.5
Default - Step by Step
51\%
-0.020 [-0.071, 0.025]
p = 0.470
Gemini 2.0 Flash
Default - Step by Step
Average
0.062 [0.028, 0.095]
p < 0.001
Gemini 2.0 Flash
Default - Step by Step
100\%
0.071 [0.005, 0.136]
p = 0.039
Gemini 2.0 Flash
Default - Step by Step
90\%
0.111 [0.045, 0.177]
p < 0.001
Gemini 2.0 Flash
Default - Step by Step
51\%
0.051 [-0.005, 0.106]
p = 0.093
GPT-4o-mini
Default - Step by Step
Average
0.004 [-0.021, 0.031]
p = 0.728
GPT-4o-mini
Default - Step by Step
100\%
-0.005 [-0.045, 0.040]
p = 0.730
GPT-4o-mini
Default - Step by Step
90\%
-0.010 [-0.051, 0.030]
p = 0.540
GPT-4o-mini
Default - Step by Step
51\%
-0.005 [-0.061, 0.051]
p = 0.788
GPT-4o
Default - Step by Step
Average
0.069 [0.023, 0.113]
p = 0.003
GPT-4o
Default - Step by Step
100\%
-0.010 [-0.061, 0.040]
p = 0.639
GPT-4o
Default - Step by Step
90\%
0.020 [-0.035, 0.076]
p = 0.431
GPT-4o
Default - Step by Step
51\%
0.061 [-0.010, 0.126]
p = 0.096
Table S5. Pairwise-comparison results between direct and CoT responses for each reasoning
model using paired bootstrap-permutation tests (5,000 replicates). P-values represent the
proportion of permuted differences with absolute values exceeding observed differences under
the null hypothesis. Questions were considered 'perfect' at ≥100 successes, with sensitivity
analyses at thresholds of 90 and 51.
Model
Conditions
Metric
RD [95\% CI]
Statistics
o3-mini
Direct - Step by Step
Average
0.029 [0.003, 0.053]
p = 0.024
o3-mini
Direct - Step by Step
100\%
0.015 [-0.056, 0.086]
p = 0.738
o3-mini
Direct - Step by Step
90\%
0.030 [-0.025, 0.086]
p = 0.332
o3-mini
Direct - Step by Step
51\%
0.000 [-0.040, 0.040]
p = 1.000
o4-mini
Direct - Step by Step
Average
0.031 [0.010, 0.052]
p = 0.003
o4-mini
Direct - Step by Step
100\%
0.010 [-0.061, 0.081]
p = 0.731
o4-mini
Direct - Step by Step
90\%
0.051 [-0.005, 0.106]
p = 0.106
o4-mini
Direct - Step by Step
51\%
0.056 [0.015, 0.096]
p = 0.003
Flash 2.5
Direct - Step by Step
Average
-0.033 [-0.056, -0.010]
p = 0.005
Flash 2.5
Direct - Step by Step
100\%
-0.131 [-0.207, -0.056]
p = 0.001
Flash 2.5
Direct - Step by Step
90\%
-0.071 [-0.131, -0.010]
p = 0.016
Flash 2.5
Direct - Step by Step
51\%
-0.035 [-0.076, 0.005]
p = 0.060
\end{document}
